% 第 5 章 实验与结果分析
%===============================================================================
% 修改记录:
%   2026-05-05  创建
%===============================================================================

\chapter{实验与结果分析}
\label{ch:experiment}

\section{实验设置}
\label{sec:setup}

实验硬件平台为基于 STM32F407VGT6 的自制开发板，主频锁定 168~MHz，FSMC 工作
于 84~MHz（HCLK/2）。SRAM 选用 IS62WV51216BLL-55TLI（55~ns 异步），LCD
为 ILI9341（240$\times$320，16 位并口）。开发环境 STM32CubeIDE 1.13 + HAL
库 V1.27，固件优化级别 \texttt{-O2}。性能测量基于 Cortex-M4 内置的 DWT 周期
计数器，单次测试取 100 次平均。对比基准包括厂家推荐配置~\cite{stm32rm0090}
与文献~\cite{zhang2021timing}所提方案。

\section{实验结果}
\label{sec:results}

\subsection{SRAM 写带宽随 DATAST 的变化}

固定 \texttt{ADDSET}=1、\texttt{ADDHLD}=1，扫描 \texttt{DATAST} 取值 1--9，
测得 SRAM 写入 1~MB 数据的实际带宽如图~\ref{fig:exp_bw}所示。当
\texttt{DATAST}$\le$3 时出现读回错误，故有效区间为 4--9。本文自适应方法
在 \texttt{DATAST}=4 处停下，达到 58.7~MB/s，较厂家推荐的 \texttt{DATAST}=9
（35.2~MB/s）提升 67\%。

\begin{figure}[htbp]
    \centering
    \begin{tikzpicture}
        \begin{axis}[thesis,
            xlabel={DATAST (HCLK 周期)}, ylabel={写带宽 (MB/s)},
            xtick={4,5,6,7,8,9}, ymin=30, ymax=65,
            legend pos=north east]
            \addplot+[mark=*] coordinates
                {(4,58.7) (5,52.1) (6,46.8) (7,42.3) (8,38.4) (9,35.2)};
            \addlegendentry{实测带宽}
            \addplot+[mark=square*, dashed] coordinates
                {(4,60.0) (5,53.3) (6,48.0) (7,43.6) (8,40.0) (9,36.9)};
            \addlegendentry{理论上限}
        \end{axis}
    \end{tikzpicture}
    \caption{不同 DATAST 配置下的 SRAM 写带宽}
    \label{fig:exp_bw}
\end{figure}

\subsection{多方案延迟对比}

将本文方法与厂家推荐配置、文献~\cite{zhang2021timing}方案及无 FSMC 的
GPIO 模拟方案在同一硬件平台上对比，LCD 全屏刷新（240$\times$320$\times$2~B
$\approx$ 150~KB）的耗时如图~\ref{fig:exp_latency}所示。

\begin{figure}[htbp]
    \centering
    \begin{tikzpicture}
        \begin{axis}[thesis,
            ybar, bar width=14pt,
            symbolic x coords={GPIO 模拟, 厂家推荐, 文献\cite{zhang2021timing}, 本文方法},
            xtick=data, x tick label style={font=\footnotesize, rotate=15, anchor=east},
            ymin=0, ylabel={LCD 全屏刷新延迟 (ms)},
            nodes near coords, nodes near coords style={font=\footnotesize}]
            \addplot coordinates {(GPIO 模拟,142) (厂家推荐,28) (文献\cite{zhang2021timing},22) (本文方法,18)};
        \end{axis}
    \end{tikzpicture}
    \caption{各方案 LCD 全屏刷新延迟对比}
    \label{fig:exp_latency}
\end{figure}

\section{分析讨论}
\label{sec:discussion}

实验结果表明，本文所提自适应搜索方法在不依赖人工查表的前提下，使 SRAM
写带宽逼近理论上限，剩余约 2\% 的差距来源于 AHB 仲裁开销与 DWT 计数器
的边界误差，符合预期。LCD 刷新延迟相比 GPIO 模拟降低 87\%，相比厂家
推荐降低 36\%，验证了 FSMC 在高速并行接口场景下的不可替代性。

需要指出的是，本文实验仅覆盖室温（25~°C）单一工况。当工作温度变化或
器件老化时，本文方法应配合一定的温度裕量（建议在搜索得到的最小值上加
1 拍）以保证长期可靠性。
